\documentclass[10pt,a4paper]{amsart}
\usepackage{amsmath,amssymb,amsthm,graphicx}
\usepackage{verbatim}
\newenvironment{reference}{\comment}{\endcomment}
%\newenvironment{reference}{}{}
\newenvironment{slogan}{\comment}{\endcomment}
\newenvironment{history}{\comment}{\endcomment}

% For commutative diagrams we use Xy-pic
\usepackage[all]{xy}

% We use 2cell for 2-commutative diagrams.
\xyoption{2cell}
% Theorem environments.
%
\theoremstyle{plain}
\newtheorem{theorem}[subsection]{Theorem}
\newtheorem{proposition}[subsection]{Proposition}
\newtheorem{lemma}[subsection]{Lemma}

\theoremstyle{definition}
\newtheorem{definition}[subsection]{Definition}
\newtheorem{example}[subsection]{Example}
\newtheorem{exercise}[subsection]{Exercise}
\newtheorem{situation}[subsection]{Situation}

\theoremstyle{remark}
\newtheorem{remark}[subsection]{Remark}
\newtheorem{remarks}[subsection]{Remarks}

\numberwithin{equation}{subsection}

% Macros
%
\def\lim{\mathop{\mathrm{lim}}\nolimits}
\def\colim{\mathop{\mathrm{colim}}\nolimits}
\def\Spec{\mathop{\mathrm{Spec}}}
\def\Hom{\mathop{\mathrm{Hom}}\nolimits}
\def\Ext{\mathop{\mathrm{Ext}}\nolimits}
\def\SheafHom{\mathop{\mathcal{H}\!\mathit{om}}\nolimits}
\def\SheafExt{\mathop{\mathcal{E}\!\mathit{xt}}\nolimits}
\def\Sch{\mathit{Sch}}
\def\Mor{\mathop{\mathrm{Mor}}\nolimits}
\def\Ob{\mathop{\mathrm{Ob}}\nolimits}
\def\Sh{\mathop{\mathit{Sh}}\nolimits}
\def\NL{\mathop{N\!L}\nolimits}
\def\CH{\mathop{\mathrm{CH}}\nolimits}
\def\proetale{{pro\text{-}\acute{e}tale}}
\def\etale{{\acute{e}tale}}
\def\QCoh{\mathit{QCoh}}
\def\Ker{\mathop{\mathrm{Ker}}}
\def\Im{\mathop{\mathrm{Im}}}
\def\Coker{\mathop{\mathrm{Coker}}}
\def\Coim{\mathop{\mathrm{Coim}}}

% Boxtimes
%
\DeclareMathSymbol{\boxtimes}{\mathbin}{AMSa}{"02}

%
% Macros for moduli stacks/spaces
%
\def\QCohstack{\mathcal{QC}\!\mathit{oh}}
\def\Cohstack{\mathcal{C}\!\mathit{oh}}
\def\Spacesstack{\mathcal{S}\!\mathit{paces}}
\def\Quotfunctor{\mathrm{Quot}}
\def\Hilbfunctor{\mathrm{Hilb}}
\def\Curvesstack{\mathcal{C}\!\mathit{urves}}
\def\Polarizedstack{\mathcal{P}\!\mathit{olarized}}
\def\Complexesstack{\mathcal{C}\!\mathit{omplexes}}
% \Pic is the operator that assigns to X its picard group, usage \Pic(X)
% \Picardstack_{X/B} denotes the Picard stack of X over B
% \Picardfunctor_{X/B} denotes the Picard functor of X over B
\def\Pic{\mathop{\mathrm{Pic}}\nolimits}
\def\Picardstack{\mathcal{P}\!\mathit{ic}}
\def\Picardfunctor{\mathrm{Pic}}
\def\Deformationcategory{\mathcal{D}\!\mathit{ef}}
\usepackage{cleveref}
\crefname{originalreference}{original reference}{original references}
\Crefname{originalreference}{Original reference}{Original references}
\makeatletter
\newlabel{lemma-characterize-normal}{{09HQ}{source}{}{sourcepointer-1}{}}
\newlabel{lemma-characterize-normal@cref}{{[originalreference][0][]09HQ}{[1][0][]source}{}{}{}}
\makeatother
\crefname{originalreference}{original reference}{original references}
\Crefname{originalreference}{Original reference}{Original references}
\makeatletter
\newlabel{lemma-map-into-algebraic-closure}{{09GU}{source}{}{sourcepointer-2}{}}
\newlabel{lemma-map-into-algebraic-closure@cref}{{[originalreference][0][]09GU}{[1][0][]source}{}{}{}}
\makeatother
\crefname{originalreference}{original reference}{original references}
\Crefname{originalreference}{Original reference}{Original references}
\makeatletter
\newlabel{theorem-existence-algebraic-closure}{{09GT}{source}{}{sourcepointer-3}{}}
\newlabel{theorem-existence-algebraic-closure@cref}{{[originalreference][0][]09GT}{[1][0][]source}{}{}{}}
\makeatother
\title{Automorphisms restrict to normal subextensions, and maps into a normal extension extend to automorphisms}
\author{The Stacks Project Authors}
\date{stacks-project:fields.tex:local-snapshot-20261012; GNU FDL 1.2 or later}
\begin{document}
\maketitle
\noindent\textit{Editorial note.} All delivered mathematical text is the author's own, copied from the pinned Stacks Project chapter file fields.tex; no statement or proof step was rewritten, reordered or removed.\par\medskip
\noindent\textit{Editorial note.} The counted proof is the author's complete argument: The xymatrix diagram (source lines 2181-2185) is part of the counted proof span and renders with the chapter preamble's xy package loading.\par\medskip
\noindent\textit{Editorial note.} The label lemma-lift-maps names the counted statement, and the proof environment immediately following it in the source is delivered verbatim as its proof.\par\medskip
\noindent\textit{Editorial note.} The counted statement is carried by the author's own lemma environment in the chapter source.\par\medskip
\noindent\textit{Allowed standard prerequisites (editorial).}\begin{itemize}
\item Same chapter: theorem-existence-algebraic-closure (delivered card L48331), lemma-map-into-algebraic-closure (delivered card L48339: K-embeddings extend along algebraic extensions), lemma-characterize-normal (normality as the equal-image criterion for F-embeddings into an algebraic closure)
\end{itemize}\medskip
\ifdefined\hypertarget\hypertarget{sourcepointer-1}{}\fi
\noindent\textit{Source reference.} Original reference 09HQ: sources/stacks-fields/pinned-original-source.tex, source lines 2096--2096. Prerequisite pointer invoked at 2179; it names a standard fact of the same chapter that the counted argument uses, and every step of the counted proof built on it is retained verbatim in the delivered body.\par\medskip
\ifdefined\hypertarget\hypertarget{sourcepointer-2}{}\fi
\noindent\textit{Source reference.} Original reference 09GU: sources/stacks-fields/pinned-original-source.tex, source lines 1240--1240. Prerequisite pointer invoked at 2184; it names a standard fact of the same chapter that the counted argument uses, and every step of the counted proof built on it is retained verbatim in the delivered body.\par\medskip
\ifdefined\hypertarget\hypertarget{sourcepointer-3}{}\fi
\noindent\textit{Source reference.} Original reference 09GT: sources/stacks-fields/pinned-original-source.tex, source lines 1184--1184. Prerequisite pointer invoked at 2175; it names a standard fact of the same chapter that the counted argument uses, and every step of the counted proof built on it is retained verbatim in the delivered body.\par\medskip
\begin{lemma}
\label{lemma-lift-maps}
Let $L/M/K$ be a tower of algebraic extensions.
\begin{enumerate}
\item If $M/K$ is normal, then any automorphism $\tau$ of $L/K$
induces an automorphism $\tau|_M : M \to M$.
\item If $L/K$ is normal, then any $K$-algebra map $\sigma : M \to L$
extends to an automorphism of $L$.
\end{enumerate}
\end{lemma}

\begin{proof}
Choose an algebraic closure $\overline{L}$ of $L$
(Theorem \ref{theorem-existence-algebraic-closure}).

\medskip\noindent
Let $\tau$ be as in (1). Then $\tau(M) = M$ as subfields of $\overline{L}$
by Lemma \ref{lemma-characterize-normal} and hence
$\tau|_M : M \to M$ is an automorphism.

\medskip\noindent
Let $\sigma : M \to L$ be as in (2).
By Lemma \ref{lemma-map-into-algebraic-closure}
we can extend $\sigma$ to a map
$\tau : L \to \overline{L}$, i.e., such that
$$
\xymatrix{
L \ar[r]_\tau & \overline{L} \\
M \ar[u] \ar[ru]_\sigma & K \ar[l] \ar[u]
}
$$
is commutative. By Lemma \ref{lemma-characterize-normal} we see that
$\tau(L) = L$. Hence $\tau : L \to L$ is an automorphism which
extends $\sigma$.
\end{proof}

\end{document}
